\section{Introduction}\label{sec:intro}

Human genomes vary across individuals at many scales, from small variants to large structural changes, yet most analysis still relies predominantly on mapping sequencing reads to a single linear reference genome. While effective in conserved regions, this paradigm breaks down in highly polymorphic loci and in sequences absent from the reference, leading to systematic reference bias~\cite{lin2024measuring}. Pangenome graphs~\cite{eizenga2020pangenome, secomandi2025pangenome} address this limitation by representing population variation explicitly within a unified structure: nodes encode DNA segments, edges capture admissible adjacencies, and paths represent individual haplotypes. By eliminating reliance on a single reference, graph-based representations substantially improve downstream analyses, including up to 34\% fewer small-variant errors and more than doubling structural-variant discovery compared to linear-reference approaches~\cite{liao2023draft}.

Recently, the Graphical Fragment Assembly (GFA) format~\cite{gfaspec} has emerged as the de facto interchange format for pangenome graphs. GFA is simple, human-readable, and widely adopted by graph construction and analysis pipelines such as Minigraph-Cactus~\cite{hickey2024pangenome} and PGGB~\cite{garrison2024building}. However, this generality of GFA comes at a cost of scale. While graph topology grows sublinearly as haplotypes are added, traversal records grow nearly linearly because each traversal is materialized explicitly. This structural redundancy rapidly dominates file size, making large GFA files expensive to store, transfer, and process, and thus motivating the need for effective file compression.

Real pangenome datasets make this bottleneck explicit. Sirén and Paten~\cite{siren2022gbz} report that 90-haplotype human pangenome graphs constructed from year-1 HPRC data require 44.9\,GiB and 88.6\,GiB as uncompressed GFA text for Cactus and PGGB, respectively; even after gzip compression, they remain 11.1\,GiB and 14.6\,GiB. At larger cohort scale, their 5008-haplotype 1000GP graph reaches 9534.9\,GiB uncompressed and 2231.3\,GiB after gzip compression. This growth is largely driven by traversal data: in a human chromosome 19 graph with 1{,}000 haplotypes, traversals alone account for 98.7\% of a 14\,GB GFA file~\cite{heringer2025human}.

In practice, gzip-compatible compressors such as gzip and bgzip remain widely used in bioinformatics pipelines because they are ubiquitous, stable, and interoperable with existing Unix tooling. However, pangenome graphs expose a mismatch between general-purpose compression and graph structure. Gzip-style compressors operate over local byte windows and are not designed to exploit repeated traversal structure across chromosome-scale haplotypes, leaving substantial cross-haplotype redundancy unexploited.

Two specialized approaches illustrate the current design trade-offs. GBZ~\cite{siren2022gbz} achieves weaker compression on highly collapsed regions and small graphs, and is not designed for incremental updates. In contrast, sqz~\cite{heringer2025human} preserves a text workflow but is relatively slow in practice due to its heap-based Byte Pair Encoding (BPE) design. Both approaches do not preserve all original GFA record information, such as optional fields, which limits their use as fully lossless GFA text compressors.

We thus present \gfaz, a GFA compressor that achieves state-of-the-art compression ratio and high throughput on both CPU and GPU. In summary, \gfaz makes four main contributions:

\begin{itemize}[noitemsep, topsep=2pt, leftmargin=*]
    \item \textbf{Orientation-aware grammar compression}: multi-round 2-mer grammar that exploits the canonical symmetry of forward/reverse node orientations with linear-time compression and decompression.
    \item \textbf{Columnar GFA representation}: column-oriented layout for all GFA record types with type-specific encodings, enabling individual field access without full decompression.
    \item \textbf{Incremental haplotype extension and random access}: support for adding extra haplotypes without recompressing previously compressed data, and for retrieving individual haplotypes without decompressing the full dataset.
    \item \textbf{Fully parallel CPU and GPU backends}: end-to-end parallel compression and decompression using OpenMP and CUDA, with a shared format so that data compressed on CPU can be decompressed on GPU for faster decoding.
\end{itemize}

Across large human pangenome datasets, \gfaz achieves 18--84$\times$ compression, typically over 10$\times$ better than gzip/zstd baselines. CPU throughput reaches 200--400 MB/s for compression and 2 GB/s for decompression; the GPU backend reaches up to 5 GB/s and 10 GB/s, respectively. These results make \gfaz practical for large-scale pangenome storage and transfer. We release our implementation as open source at: \href{https://github.com/babyplutokurt/GFAz}{https://github.com/babyplutokurt/GFAz}.
